% Generated draft: requires mathematical and source review.
\documentclass[11pt]{article}
\usepackage[margin=25mm]{geometry}
\usepackage{fontspec}
\setmainfont{Latin Modern Roman}
\usepackage{amsmath,amssymb,mathtools,unicode-math}
\setmathfont{Latin Modern Math}
\usepackage{enumitem,needspace,xurl,hyperref}
\setlength{\parindent}{0pt}
\setlength{\parskip}{5pt}
\setlength{\emergencystretch}{3em}
\newcommand{\N}{\mathbb{N}}
\newcommand{\Z}{\mathbb{Z}}
\newcommand{\Q}{\mathbb{Q}}
\newcommand{\R}{\mathbb{R}}
\newcommand{\C}{\mathbb{C}}
\newcommand{\F}{\mathbb{F}}
\newcommand{\A}{\mathbb{A}}
\newcommand{\PP}{\mathbb{P}}
\newcommand{\E}{\mathbb{E}}
\newcommand{\eps}{\varepsilon}
\newcommand{\dd}{\,\mathrm{d}}
\begin{document}
\section*{UnsolvedMath ID: 700257}
\textbf{Source problem code:} catalogue entry
\subsection{Definitions and mathematical statement}
A variety of groups is a class \(\mathcal V\) consisting of all groups satisfying a specified collection of group identities. For a group \(G\), set \(G^{(0)}=G\) and \(G^{(n+1)}=[G^{(n)},G^{(n)}]\), where \([K,K]\) denotes the subgroup generated by the commutators \(x^{-1}y^{-1}xy\) with \(x,y\in K\). A soluble variety is understood here in the standard sense that there exists an integer \(d\geq 0\) such that \(G^{(d)}=\{1\}\) for every \(G\in\mathcal V\).

A group \(H\) is residually finite if, for every \(h\in H\setminus\{1\}\), there exist a finite group \(F\) and a homomorphism \(\varphi:H\to F\) with \(\varphi(h)\neq 1\).

\textbf{Classification problem.} Determine all soluble varieties \(\mathcal V\) such that
\[
\forall G\in\mathcal V\;\forall H\leq G,\qquad
H\text{ finitely generated}\ \Longrightarrow\ H\text{ residually finite}.
\]
Since a variety is closed under taking subgroups, this is equivalently the requirement that every finitely generated group belonging to \(\mathcal V\) be residually finite.

The source phrase ``all of whose finitely generated subgroups'' is interpreted as quantifying over the finitely generated subgroups of every group in the variety. The supplied record does not explicitly expand this quantification. It requests a classification without specifying a form for that classification.
\subsection{Short English statement}
Classify the soluble varieties of groups for which every finitely generated subgroup of every member is residually finite. Residual finiteness means that each nonidentity element remains nonidentity under some homomorphism to a finite group. This expands the source's implicit quantification over members of a variety.
\subsection{Sources}
\begin{itemize}
\item Local UnsolvedMath registry, permanent ID \(700257\), exact source code \texttt{catalogue entry}, title \emph{Kourovka Notebook Problem 4.50}. The supplied record labels the problem \texttt{open}; that label is supplied metadata, not an independently verified mathematical or literature assessment.
\item \emph{The Kourovka Notebook (Unsolved Problems in Group Theory)}, original Problem 4.50, fourth issue (1973); proposed by V. N. Remeslennikov.
\item The supplied background identifies the twenty-first issue edition (2026), \texttt{arXiv:1401.0300v48}. The collection description names E. I. Khukhro and V. D. Mazurov as editors. These bibliographic details are reproduced from the supplied record; the cited edition has not been independently consulted.
\end{itemize}
\end{document}
